\newpage
\section{Acta Número 27}
\label{sec:acta27}

\noindent \textbf{Fecha de la sesión:}\par
Jueves, 04 de Junio de 2026

\vspace*{0.5cm}
\noindent \textbf{Datos de la Sesión:}
\vspace{-0.2cm}
\begin{itemize}[noitemsep]
    \item \textbf{Modalidad:} Virtual - Google Meet.
    \item \textbf{Hora de Inicio:} 18:00
    \item \textbf{Hora de Fin:} 20:00
\end{itemize}

\vspace*{0.5cm}
\noindent \textbf{Orden del Día:}
\vspace{-0.2cm}
\begin{itemize}[noitemsep]
    \item Control de Asistencia.
    \item Ceremonia de Entrega de Documentación del Sprint.
    \item Ceremonia de Finalización e Incremento (Revisión DoD, Integración, Deployment y Smoke Test).
    \item Ceremonia de Sprint Review y Retrospectiva del ciclo.
    \item Revisión final de entregables, resolución de dudas y cierre del ciclo.
    \item Firmas y Aprobación de los Presentes.
\end{itemize}

\vspace*{0.5cm}
\noindent \textbf{Socios Fundadores Presentes:}
\vspace{-0.2cm}
\begin{enumerate}[noitemsep]
    \item Pardo Pozo Juan Diego (Jefe de Proyecto).
    \item Arnez Jaimes Mauricio.
    \item Quispe Flores Camila Belen.
    \item Ariñez Bautista Jim Gabriel.
    \item Medina Rodríguez Mateo Jhoustin.
    \item Molina Maldonado Tomas Manuel.
    \item Gutierrez Guzmán Angheline.
\end{enumerate}

\newpage
\subsection{Desarrollo del Acta}

\noindent \textbf{Control de Asistencia del Team}\par
Se procedió a la toma de asistencia a la hora acordada, corroborando la presencia de los 7 socios fundadores de la grupo-empresa Byte Busters S.R.L. Al verificar el quórum correspondiente y la completa atención del equipo, se dio inicio formal a esta sesión de cierre y validación.

\vspace*{0.5cm}
\noindent \textbf{Ceremonia de Entrega de Documentación del Sprint}\par
El primer punto de trabajo consistió en la consolidación de la etapa documental. El equipo realizó la entrega formal de la documentación del producto correspondiente a este sprint, verificando que los manuales, especificaciones técnicas y reportes de QA reflejen con total exactitud las características del software desarrollado. Se destacó el esfuerzo conjunto para mantener los documentos actualizados, asegurando que la transferencia de conocimiento cumpla con los máximos estándares de calidad.

\vspace*{0.5cm}
\noindent \textbf{Ceremonia de Finalización e Incremento}\par
Inmediatamente después, se ejecutaron las actividades críticas para la liberación del producto. En primer lugar, se llevó a cabo la Revisión Final de la Definición de Hecho (DoD), confirmando el cumplimiento absoluto de las normativas de calidad. Superado este filtro, se procedió con la Integración y Build, generando el incremento funcional de forma exitosa. Acto seguido, el equipo ejecutó la Puesta en Staging (Deployment) sobre los entornos de \textbf{Vercel} y \textbf{Supabase} y finalizó el proceso aplicando un Smoke Test Final, el cual certificó la estabilidad, robustez y correcta finalización del incremento.

\vspace*{0.5cm}
\noindent \textbf{Ceremonia de Sprint Review y Retrospectiva}\par
Con el incremento validado, el equipo llevó a cabo la Sprint Review, presentando las funcionalidades terminadas y contrastándolas con el Sprint Goal definido al inicio del ciclo, dando por aceptados los Product Backlog Items completados. A continuación, se desarrolló la Retrospectiva, en la que se analizaron de manera honesta los aciertos del ciclo, las dificultades encontradas durante la migración a los entornos externos y las oportunidades de mejora. Se registraron las acciones correctivas para incorporarlas en la planificación del siguiente sprint.

\vspace*{0.5cm}
\noindent \textbf{Revisión Final de Entregables y Cierre}\par
Para concluir, se realizó una última revisión panorámica de todos los entregables del ciclo. Al abrir un espacio para consultas y observaciones finales, se evidenció una claridad total por parte del equipo técnico, confirmando que no existían dudas pendientes respecto a la arquitectura, documentación o pasos a seguir. Habiendo completado exitosamente las ceremonias de documentación, incremento, review y retrospectiva, el equipo manifiesta su total y unánime conformidad con los resultados del sprint. Se da por aprobada la sesión y se procede con la firma de los presentes.

\newpage
\noindent \textbf{Firmas y Aprobación de los Presentes}
\begin{center}
    \noindent
    \setlength{\tabcolsep}{0pt}
    \begin{tabular}{ccc}
        \espacioFirma{assets/img/firmas/mauricio.jpg}{Mauricio Jaimes Arnez}{13751221} &
        \espacioFirma{assets/img/firmas/camila.jpg}{Camila Belen Quispe Flores}{13097244} &
        \espacioFirma{assets/img/firmas/jim.jpg}{Jim Gabriel Ariñez Bautista}{9517642} \\[1.0cm]
        \espacioFirma{assets/img/firmas/mateo.jpg}{Mateo Medina Rodríguez}{9453809} &
        \espacioFirma{assets/img/firmas/tomas.jpg}{Tomas Molina Maldonado}{8051701} &
        \espacioFirma{assets/img/firmas/angheline.jpg}{Angheline Gutierrez Guzmán}{13751815} \\[1.0cm]
    \end{tabular}
    \vspace{0.2cm}
    \begin{tabular}{cc}
        \espacioFirma{assets/img/firmas/diego.jpg}{Juan Diego Pardo Pozo}{12747783}
    \end{tabular}
\end{center}

\vspace{0.5cm}
\noindent \textbf{Fecha de última revisión:} Jueves, 04 de Junio de 2026
